% Assignment 5 (Q2) - Beamer presentation supporting the report
% Tarun Konde (2025UCP1811)

\documentclass{beamer}

\usetheme{Madrid}
\usecolortheme{default}

\title[Introduction to \LaTeX]{Introduction to \LaTeX}
\subtitle{Technical Writing Lab -- Assignment 5}
\author[Tarun Konde]{Tarun Konde \\ 2025UCP1811}
\institute[MNIT Jaipur]{Department of CSE, MNIT Jaipur}
\date{\today}

\begin{document}

% title slide
\begin{frame}
    \titlepage
\end{frame}

\begin{frame}{Outline}
    \tableofcontents
\end{frame}

\section{What is \LaTeX}

\begin{frame}{What is \LaTeX?}
    \begin{itemize}
        \item<1-> A document preparation system built on \TeX
        \item<2-> Content is written as plain text with commands
        \item<3-> The compiler handles layout, numbering and references
    \end{itemize}

    \only<4>{
    \begin{block}{Compile command}
        \texttt{pdflatex report.tex}
    \end{block}
    }
\end{frame}

\section{Features}

\begin{frame}{Math and Structure}
    \begin{columns}
        \begin{column}{0.5\textwidth}
            \textbf{Structure}
            \begin{itemize}
                \item \texttt{\textbackslash chapter}, \texttt{\textbackslash section}
                \item Automatic table of contents
                \item Headers with \texttt{fancyhdr}
            \end{itemize}
        \end{column}
        \begin{column}{0.5\textwidth}
            \textbf{Math}
            \[
                \int_0^1 x^2\,dx = \frac{1}{3}
            \]
            \[
                E = mc^2
            \]
        \end{column}
    \end{columns}
\end{frame}

\section{Conclusion}

\begin{frame}{Conclusion}
    \begin{alertblock}{Takeaway}
        \LaTeX{} has a learning curve, but it is the better tool for long
        technical documents with equations and references.
    \end{alertblock}

    \vspace{0.5cm}
    \centering
    Thank you
\end{frame}

\end{document}
